\chapter{Dérivation alternative des fluctuations de $\rho$}\label{annex:Jerome}

Dans cette annexe, nous proposons une seconde dérivation des fluctuations de la densité de rapidités $\rho$, différente de celle présentée au chapitre~\ref{chap:Fluctu}.  
Cette approche m’a été suggérée par Jérôme Dubail (\gls{CESQ}). Elle présente l’intérêt de passer directement par les fluctuations du facteur d’occupation $\nu$, et constitue également une bonne occasion de manipuler les dérivées de l’opérateur \emph{d’habillage}.

\section{Réécriture de l’entropie de Yang--Yang}

L’entropie de Yang--Yang s’écrit
\begin{eqnarray}
	\mathcal{S}_{YY} 
	&=& \int d \theta \, \Big( \rho_s \ln \rho_s - \rho \ln \rho - (\rho_s - \rho ) \ln (\rho_s - \rho ) \Big)(\theta).    
\end{eqnarray}

En introduisant le facteur d’occupation
\begin{eqnarray}
	\nu = \frac{\rho}{\rho_s},	
\end{eqnarray}
l’intégrande peut se réécrire sous la forme
\begin{eqnarray}
	\mathcal{S}_{YY} 
	&=& \int d\theta \, s(\nu(\theta))\, \rho_s(\theta), 	
\end{eqnarray}
où la fonction d’entropie locale est donnée par
\begin{eqnarray}
	s(\nu) &=& - \nu \ln \nu - (1 - \nu ) \ln (1 - \nu ).
\end{eqnarray}

\section{Différentielle de l’action effective}

On définit l’action effective comme
\begin{eqnarray}
	\mathcal{S}_{YY} - \mathcal{W},
\end{eqnarray}
où $\mathcal{W}$ est l’énergie généralisée,
\begin{eqnarray}
	\mathcal{W} &=& \int d\theta \, w(\theta)\, \rho(\theta).		
\end{eqnarray}

Sa différentielle est donnée par
\begin{eqnarray}\label{annex.J.1}
	\delta ( \mathcal{S}_{YY} -  \mathcal{W} ) 
	&=& \int d \theta \, \Big( \delta \nu \, s'(\nu)\, \rho_s + s(\nu)\, \delta \rho_s - w\, \delta \rho \Big)(\theta).
\end{eqnarray}
En remarquant que
\begin{eqnarray}\label{annex.J.rho}
	\delta \nu \, \rho_s = \delta \rho - \nu \, \delta \rho_s, 	
\end{eqnarray}
on réécrit~\eqref{annex.J.1} sous la forme
\begin{eqnarray}
	\delta ( \mathcal{S}_{YY} -  \mathcal{W} ) 
	&=& \int d \theta \, \Big( \{  s'(\nu) - w\} \delta \rho  
	+  \{  s(\nu) - \nu \, s'(\nu) \} \delta \rho_s \Big)(\theta). 		
\end{eqnarray}

On obtient alors la différentielle seconde :
\begin{eqnarray}
	\delta^2 ( \mathcal{S}_{YY} -  \mathcal{W} )  
	&=& \int d \theta \, \Big( \delta \nu \, s''(\nu)\, \delta \rho  
	- \nu \, \delta \nu \, s''(\nu)\, \delta \rho_s \Big)(\theta).
\end{eqnarray}

En utilisant~\eqref{annex.J.rho}, on simplifie pour obtenir
\begin{eqnarray}\label{annex.J.2}
	\delta^2 ( \mathcal{S}_{YY} -  \mathcal{W} )  
	&=& \int d \theta \, \big( (\delta \nu )^2 \, s''(\nu)\, \rho_s \big)(\theta),		
\end{eqnarray}
avec
\begin{eqnarray}
	s''(\nu) &=& - \frac{1}{\nu (1 - \nu)}.	
\end{eqnarray}

\section{Fluctuations}

\subsection{Fluctuations des facteurs d’occupation}

Les fluctuations du facteur d’occupation s’écrivent
\begin{eqnarray}
	\braket{\delta\nu(\theta')  \, \delta\nu(\theta)}_w	
	&=& - \Bigg[ L \frac{ \delta^2 ( \mathcal{S}_{YY} -  \mathcal{W} )}{\delta \nu \, \delta \nu } \Bigg]^{-1} ( \theta , \theta' ).
\end{eqnarray}
En injectant~\eqref{annex.J.2}, il vient
\begin{eqnarray}
	\braket{\delta\nu(\theta')  \, \delta\nu(\theta)}_w	
	&=& -\, s''(\nu(\theta)) \, \rho_s(\theta)\, \delta ( \theta - \theta' ),		
\end{eqnarray}
qui est purement diagonale.

\subsection{Opérateur d’habillage}

On rappelle la définition de l’opérateur d’habillage :
\begin{eqnarray}
	f^{\mathrm{dr}}_{[\nu]} = f + \frac{\Delta}{2\pi} \star ( \nu \, f^{\mathrm{dr}}_{[\nu]} ).	
\end{eqnarray}
En remarquant que 
\begin{eqnarray}
	C(\theta , \lambda) = \Bigg[ \frac{\Delta ( \lambda - \cdot ) }{2\pi } \Bigg]^{\mathrm{dr}}_{[\nu]} (\theta), 	
\end{eqnarray}
qui est symétrique, on peut réécrire~\cite{Doyon_2017}
\begin{eqnarray}
	f^{\mathrm{dr}}_{[\nu]}(\theta) 
	&=&  \int d\lambda \, f(\lambda) \,  \Big( \delta ( \theta - \lambda ) + \nu(\lambda)\, C(\theta,\lambda) \Big).
\end{eqnarray}

\subsection{Fluctuations des distributions de rapidité}

On a vu que
\begin{eqnarray}
	2\pi \, \rho_s (\theta) 
	&=& 1^{\mathrm{dr}}_{[\nu]} (\theta) 
	= \int d \lambda \, \Big( \delta ( \theta - \lambda ) + \nu(\theta)\, C(\theta,\lambda) \Big),
\end{eqnarray}
d’où
\begin{eqnarray}
	2\pi \, \delta \rho_s (\theta) 
	&=& \int d \lambda \, C(\theta,\lambda)\, \rho_s(\lambda)\, \delta \nu(\lambda).	
\end{eqnarray}

Ainsi,
\begin{eqnarray}
	\delta \rho (\theta) 
	&=& \rho_s(\theta) \delta \nu (\theta) + \nu(\theta) \delta \rho_s (\theta) \\
	&=& \int d \lambda \, \delta ( \theta - \lambda ) \rho_s(\lambda)\, \delta \nu(\lambda)  
	+ \int d \lambda \, \nu (\theta)\, C(\theta,\lambda)\,\rho_s(\lambda)\, \delta \nu (\lambda),  \\
	&=& \int d \lambda \, \Big( \delta ( \theta - \lambda ) + \nu (\theta)\, C(\theta,\lambda) \Big) \rho_s(\lambda)\, \delta \nu (\lambda). \label{annex.J.rho.2}
\end{eqnarray}

En réinjectant~\eqref{annex.J.2} et~\eqref{annex.J.rho.2} dans la définition
\begin{eqnarray}
	\braket{\delta\rho(\theta')  \, \delta\rho(\theta)}_w	
	&=& -  \Bigg[ L \frac{ \delta^2 ( \mathcal{S}_{YY} -  \mathcal{W} )}{\delta \rho \, \delta \rho }\Bigg]^{-1} ( \theta , \theta'),		
\end{eqnarray}
on obtient finalement
\begin{eqnarray}\label{annex.J.3}
	\braket{\delta\rho(\theta')  \, \delta\rho(\theta)}_w  
	&=& -\frac{1}{L}  \int d \lambda \, \frac{ \Big( \delta ( \theta - \lambda ) + \nu (\theta)\, C(\theta,\lambda) \Big) \Big( \delta ( \theta' - \lambda ) + \nu (\theta')\, C(\theta',\lambda) \Big)}{s''(\nu(\lambda))} \, \rho_s (\lambda) \\
	&=& \frac{1}{L}\, \mathcal{D}^{-1}(\theta , \theta') + \frac{1}{L}\, \mathcal{B}(\theta , \theta'),		
\end{eqnarray}
où
\begin{eqnarray}
	\mathcal{D}^{-1}(\theta , \theta' ) &=& \big( \rho_s(\theta)\, \nu(\theta)\, (1 - \nu(\theta)) \big)\, \delta ( \theta - \theta' ), \\
	\mathcal{B}(\theta , \theta' ) 
	&=& \Big[ \nu (\theta')\, \rho_s(\theta)\, \nu(\theta)(1 - \nu(\theta)) 
	+ \nu (\theta)\, \rho_s(\theta')\, \nu(\theta')(1 - \nu(\theta')) \Big]\, C(\theta,\theta') \\
	&& + \nu (\theta)\, \nu (\theta') \int d \lambda \, \rho_s(\lambda)\, \nu(\lambda)\,(1 - \nu(\lambda))\, C(\theta,\lambda)\, C(\theta',\lambda).
\end{eqnarray}

Cette expression coïncide numériquement avec le résultat obtenu en~\eqref{chap:fluctu:eq:fluctuations.10}.



















